
\section{Chiral Hamiltonian reduction} \label{section:brst}

We recall the BRST formalism for chiral Hamiltonian reduction. Our main references are \cite{feigin1990quantization,kac2003quantum,arakawa2007representation,arakawa2015localization}. See also \cite[Section 3.5]{genra2025reduction} and references therein fore more details on the construction of the complex and the geometric interpretation. 


\subsection{The fermionic vertex algebra}

Let $\V(\g_{\geqslant 1})$ be the vertex subalgebra of $\V^k(\g)$ that is generated by the subspace $\g_{\geqslant 1}$ of $\g$. It is a graded subalgebra for the grading \eqref{equation:grading-universal-affine}.


Denote by $\F(\g_{\geqslant 1} \oplus (\g_{\geqslant 1})^\vee)$ the \emph{fermionic vertex superalgebra} associated to the vector space $\g_{\geqslant 1} \oplus (\g_{\geqslant 1})^\vee$ equipped with the duality pairing \cite[Section 1.8]{desole2006finite}. It is a vertex algebraic analogue of a Clifford algebra, see \cite[Section 2.4.2]{genra2025reduction}. It contains $\g_{\geqslant 1} \oplus (\g_{\geqslant 1})^\vee$ as a linear subspace, and it is freely generated by any basis of $\g_{\geqslant 1} \oplus (\g_{\geqslant 1})^\vee$. If $a$ is an element of $\g_{\geqslant 1} \oplus (\g_{\geqslant 1})^\vee$, denote by $\phi(a)$ its image in the vertex superalgebra.

There is a vertex algebra grading
\[ \F(\g_{\geqslant 1} \oplus (\g_{\geqslant 1})^\vee) = \bigoplus_{\Delta \in \Z} \F(\g_{\geqslant 1} \oplus (\g_{\geqslant 1})^\vee)(\Delta) \] 
such that $\g_i$ is contained in $\F(\g_{\geqslant 1} \oplus (\g_{\geqslant 1})^\vee)(1 - i)$ and $(\g_i)^\vee$ is contained in $\F(\g_{\geqslant 1} \oplus (\g_{\geqslant 1})^\vee)(i)$ \cite[Equation (5.3)]{desole2006finite}. In the particular case of this good grading for minimal $f$, $\g_{\geqslant 1} = \g_1$ so we get a grading supported on $\Z_{\geqslant 0}$.

The vertex algebra $\F(\g_{\geqslant 1} \oplus (\g_{\geqslant 1})^\vee)$ can be equipped with the \emph{charge grading},
\[ \F(\g_{\geqslant 1} \oplus (\g_{\geqslant 1})^\vee) = \bigoplus_{n \in \Z} \F(\g_{\geqslant 1} \oplus (\g_{\geqslant 1})^\vee)^{\mathrm{ch} = n}, \] 
which is not a vertex algebra grading. Generators from $\g_{\geqslant 1}$ are of charge $-1$, generators from $(\g_{\geqslant 1})^\vee$ are of charge $+1$ and the translation operator preserves the charge.

\subsection{BRST cochain complex}

Let $\{a_i\}_{i=1}^{d-1}$ be a basis of $\g_{\geqslant 1}$ and $\{a_i^\vee\}_{i=1}^{d-1}$ be the dual basis for the duality pairing with $(\g_{\geqslant 1})^\vee$. The vertex superalgebra
\[ \V(\g_{\geqslant 1}) \otimes \F(\g_{\geqslant 1} \oplus (\g_{\geqslant 1})^\vee) \]
contains an element of charge $1$ and conformal degree $1$, called the \emph{BRST charge}:
\[ Q \coloneqq \sum_{i=1}^{d-1} a_i \otimes \phi(a_i^\vee) - \frac{1}{2} \sum_{i, j=1}^{d - 1} 1 \otimes \phi(a_i^\vee){}_{(-1)}\big(\phi(a_j^\vee){}_{(-1)}\phi([a_i, a_j])\big). \]
The zero-mode $Q_{(0)}$ is an endomorphism that squares to $0$, increases the charge-degree by $1$ and preserves the vertex algebra grading \cite[Sections 2.1 and 2.2]{kac2003quantum}. Then, the pair
\[ (\V(\g_{\geqslant 1}) \otimes \F(\g_{\geqslant 1} \oplus (\g_{\geqslant 1})^\vee), Q_{(0)}) \]
forms a cochain complex, called \emph{BRST cochain complex}, that is a graded vertex algebra. Each graded piece is a cochain subcomplex, so we have a vertex algebra grading.

Let $\widehat{\mu}_V \colon \V(\g_{\geqslant 1}) \rightarrow V$ be a (purely even) vertex algebra homomorphism, that is usually called a \emph{chiral comoment map}. We assume that $V$ is a graded vertex algebra and that $\widehat{\mu}_V$ is compatible with the gradings. The pair
\[ \big(\Cpx(V) \coloneqq V \otimes \F(\g_{\geqslant 1} \oplus (\g_{\geqslant 1})^\vee), \widehat{\mu}_V(Q)_{(0)}\big) \]
is a cochain complex that is a graded vertex algebra. Each graded piece is a cochain subcomplex. In the following proposition, denote by $[c]$ the cohomology class of some closed element $c$ in $\Cpx(V)$.

The following facts are standard, see for example \cite[Section 3.15]{arakawa2007representation}.

\begin{proposition} \label{proposition:brst-cohomology}
	The cohomology $\Hgy(\Cpx(V))$ enjoys the following properties:
	\begin{enumerate}
		\item it is a vertex superalgebra whose vacuum is $[1]$, translation operator is induced by the translation operator on the BRST complex and the state-field correspondence is given by:
		\[ [a](z)[b] \coloneqq [a(z)b] = \sum_{n \in \Z} [a_{(n)}b] \, z^{- n - 1} \quad \text{for} \quad [a], [b] \in \Hgy(\Cpx(V)),  \]
		
		\item it is a graded vertex algebra for the grading induced by the grading on the complex:
		\[ \Hgy(\Cpx(V)) = \bigoplus_{\Delta \in \Z_{\geqslant 0}}  \Hgy(\Cpx(V)(\Delta)), \]
		
		\item the degree $0$ cohomology $\Hgy^0(\Cpx(V))$ is a purely even vertex subalgebra.
	\end{enumerate}
\end{proposition}

This construction is functorial. Fix $\widehat{\mu}_V \colon \V(\g_{\geqslant 1}) \rightarrow V$ and $\widehat{\mu}_W \colon \V(\g_{\geqslant 1}) \rightarrow W$  two chiral comoment maps, and $\psi \colon V \rightarrow W$ a vertex algebra homomorphism such that $\widehat{\mu}_W = \psi \circ \widehat{\mu}_V$. Then there is a natural vertex super algebra homomorphism between the corresponding cochain complexes that descends to the corresponding cohomologies. 

\subsection{Li filtration on the quantum Hamiltonian reduction}

Let $V$ be a vertex algebra equipped with a chiral comoment map $\widehat{\mu}_V \colon \V(\g_{\geqslant 1}) \rightarrow V$. Let $X$ be an affine Poisson scheme with an action of $G_{\geqslant 1}$ and a moment map $\mu_X\colon X \rightarrow (\g_{\geqslant 1})^\vee$. We assume that $V$ is a \emph{chiral quantisation} of $X$ in the following sense: there is a vertex Poisson algebra isomorphism $\gr V \cong \C[\J_\infty X]$ that makes the triangle
\[ \begin{tikzcd}[column sep = small]
	\gr V \arrow[rr, "\sim"] &  & \C[\J_\infty X] \\
	& \arrow[lu, "\gr \widetilde{\mu}_V"] \arrow[ru, "(\J_\infty \mu_X)^*"'] \gr V(\g_{\geqslant 1})\cong \C[\J_\infty (\g_{\geqslant})^\vee]  &
\end{tikzcd} \]
commute.

Like any graded vertex superalgebra, the BRST complex $\Cpx(V)$ carries an increasing Li filtration. Because the differential preserves the vertex algebra grading and is a derivative for the modes of the vertex algebra, $\Cpx(V)$ is a filtered complex in the sense of \cite[Section 5.4]{weibel1994introduction} and $\gr \Cpx(V)$ has a natural structure of cochain complex.

There is an induced filtration on the cohomology vertex superalgebra $\Hgy(\Cpx(V))$ \cite[Section 5.5]{weibel1994introduction}: for $p$ in $\Z$, we denote by $\Hgy(\Cpx(V))_{(p)}$ the image of the homomorphism $\Hgy(\Cpx(V)_p) \rightarrow \Hgy(\Cpx(V))$. This defines an exhaustive filtration on $\Hgy(\Cpx(V))$ that is supported on $\Z_{\geqslant 0}$. 

A priori, this induced filtration is not the increasing Li filtration on the vertex superalgebra $\Hgy(\Cpx(V))$. However, this filtration inherits the properties \eqref{equation:li-filtration-property} from the Li filtration on the complex and we have the following lemma.

\begin{lemma}
	The associated graded vector space $\gr \Hgy(\Cpx(V))$ inherits a structure of vertex Poisson superalgebra.
\end{lemma}

Moreover, because both the induced filtration satisfies \eqref{equation:li-filtration-property}, and $\Hgy(\Cpx(V))_{(p)}$ contains $\Hgy(\Cpx(V)(p)) = \Hgy(\Cpx(V))(p)$ for all $p$ in $\Z$, so by \cite[Theorem 4.14]{li2004vertex}, there is an inclusion 
\[ \Hgy(\Cpx(V))_{p} \subseteq \Hgy(\Cpx(V))_{(p)}. \] 


%
%Because the grading and the vertex algebra structure on the cohomology is induced by the grading and the structure on the complex, a direct computation shows that the induced filtration is the increasing Li filtration of $\Hgy(\Cpx(V))$.

The following statement gives a sufficient condition on $X$ to compute $\gr \Hgy(\Cpx(V))$. Variations of this statement appear frequently in the literature \cite{arakawa2015associated,arakawa2015localization,kuwabara2021,,arakawa2024arc,genra2025reduction, arakawa2026faithful}. We give the proof in our particular setting.

\begin{proposition} \label{proposition:reduction-quantisation}
	We make the following technical assumption:
	\begin{enumerate}
		\item the moment map $\mu_x\colon X \rightarrow (\g_{\geqslant 1})^\vee$ is a smooth scheme homomorphism,
		
		\item there is an affine subscheme $S$ of $(\mu_X)^{-1}(0)$ such that the action of $G_{\geqslant 1}$ induces a scheme isomorphism
		\[ G_{\geqslant 1} \times S \cong (\mu_X)^{-1}(0). \]
	\end{enumerate}
	Then, the following properties hold:
	\begin{enumerate}
		\item there are scheme isomorphisms
		\[ \J_\infty S \cong \J_\infty (\mu_X)^{-1}(0) \git \J_\infty G_{\geqslant 1} \cong \J_\infty ((\mu_X)^{-1}(0) \git G_{\geqslant 1}), \]
		
		\item there is the following cohomology vanishing:
		\[ \Hgy^n(\Cpx(V)) = 0 \quad \text{for} \quad n \neq 0, \]
		
		\item the natural homomorphism
		\[ \gr \Hgy^0(\Cpx(V)) \longrightarrow \Hgy^0(\gr \Cpx(V))  \]
		is a vertex Poisson algebra isomorphism. Moreover, both are explicitly given by the arc space of the Hamiltonian reduction of $X$:
		\[ \gr \Hgy^0(\Cpx(V)) \cong \Hgy^0(\gr \Cpx(V)) \cong \J_\infty \C[(\mu_X)^{-1}(0) \git G_{\geqslant 1}].  \]
	\end{enumerate}
\end{proposition}

\begin{proof}
	The first statement is given by applying the functor $\J_\infty$ to the isomorphism $G_{\geqslant 1} \times S \cong (\mu_X)^{-1}(0)$.
	
	The increasing Li filtration on $\Cpx(V)$ is bounded from below and exhaustive. By \cite[Theorem 5.5.1]{weibel1994introduction}, it gives rise to a convergent spectral sequence $\{E_n\}_{n \geqslant 0}$ whose first pages and the infinity page are:
	\[ \begin{aligned}
		E_0^{p, q} & = \gr_p \Cpx(V)^{\mathrm{ch} = p + q}, \\
		E_1^{p, q} & = \Hgy^{p + q}(\gr_p \Cpx(V)), \\
		E_\infty^{p, q} & = \gr_p \Hgy^{P + q}(\Cpx(V)),
	\end{aligned} \]
	for $p$ and $q$ in $\Z$.
	
	The $0$-page is a vertex Poisson analogue of the BRST complex:
	\[ \gr \Cpx(V) = \C[\J_\infty X] \otimes \J_\infty \operatorname{Alt}(\g_{\geqslant 1} \oplus (\g_{\geqslant 1})^\vee). \]
	See \cite[Section 2.3.3]{arakawa2015localization} for precise definitions.
	
	Because $\mu_X$ is assumed to be smooth, the induced map $\J_\infty \mu_X \colon \J_\infty X \rightarrow \J_\infty (\g_{\geqslant 1})^\vee$ is flat by \cite[Remark 3.4.2]{genra2025reduction}. Moreover, we have assumed the isomorphism $G_{\geqslant 1} \times S \cong (\mu_X)^{-1}(0)$. Therefore, we can apply \cite[Theorem 2.3.3.1]{arakawa2015localization} and we get that
	\[ \Hgy^{n}(\gr\Cpx(V)) = 0 \quad \text{for} \quad n \neq 0. \]
	
	Therefore the $1$-page $E_1^{p, q}$ is zero outside the diagonal $p = - q$. Therefore, we have isomorphisms
	\[ E_r \cong E_{r+1} \cong E_\infty \quad \text{for} \quad r \geqslant 1, \]
	and in particular
	\[ \gr \Hgy^n(\Cpx(V)) \cong  \Hgy^n(\gr \Cpx(V)) \]
	for all $n$. But these two cohomologies are zero for $n \neq 0$. This gives the vanishing of $\Hgy^n(\Cpx(V))$ (because the filtration is bounded from below and exhaustive) and the vertex Poisson isomorphism
	\[ \gr \Hgy^0(\Cpx(V)) \cong \Hgy^0(\gr \Cpx(V)). \]
\end{proof}
